\lyseis{7}{ΜΚΔ και ΕΚΠ: μοιράζω και συναντιέμαι}
\sol{1} α) Διαιρέτες του 12: 1, 2, 3, 4, 6, 12. Του 30: 1, 2, 3, 5, 6, 10, 15, 30. Κοινοί: 1, 2, 3, 6. $\MKD=6$. \quad β) Του 18: 1, 2, 3, 6, 9, 18. Του 24: 1, 2, 3, 4, 6, 8, 12, 24. $\MKD=6$. \quad γ) Του 20: 1, 2, 4, 5, 10, 20. Του 45: 1, 3, 5, 9, 15, 45. $\MKD=5$. \quad δ) Του 16: 1, 2, 4, 8, 16. Του 27: 1, 3, 9, 27. Κοινός μόνο ο 1: $\MKD=1$.
\sol{2} α) $\EKP(4,\,6)=12$ \quad β) $\EKP(6,\,15)=30$ \quad γ) $\EKP(8,\,12)=24$ \quad δ) $\EKP(5,\,7)=35$.
\sol{3} α) $24=2^3\cdot 3$, $36=2^2\cdot 3^2$: $\MKD=2^2\cdot 3=12$, $\EKP=2^3\cdot 3^2=72$. \quad β) $60=2^2\cdot 3\cdot 5$, $84=2^2\cdot 3\cdot 7$: $\MKD=2^2\cdot 3=12$, $\EKP=2^2\cdot 3\cdot 5\cdot 7=420$. \quad γ) $90=2\cdot 3^2\cdot 5$, $150=2\cdot 3\cdot 5^2$: $\MKD=2\cdot 3\cdot 5=30$, $\EKP=2\cdot 3^2\cdot 5^2=450$. \quad δ) $48=2^4\cdot 3$, $180=2^2\cdot 3^2\cdot 5$: $\MKD=2^2\cdot 3=12$, $\EKP=2^4\cdot 3^2\cdot 5=720$.
\sol{4} $12=2^2\cdot 3$, $18=2\cdot 3^2$, $30=2\cdot 3\cdot 5$. $\MKD=2\cdot 3=6$ (κοινοί παράγοντες 2 και 3, με μικρότερο εκθέτη 1). $\EKP=2^2\cdot 3^2\cdot 5=180$.
\sol{5} Πρώτοι μεταξύ τους είναι τα (α), (β), (δ) και (στ). Τα (γ) 21 και 35 έχουν κοινό διαιρέτη το 7, και τα (ε) 14 και 49 επίσης το 7.
\sol{6} α) $\MKD=7$, $\EKP=42$ \quad β) $15$ και $60$ \quad γ) $25$ και $100$. Όταν ο ένας αριθμός διαιρεί τον άλλον, το ΜΚΔ είναι ο μικρότερος και το ΕΚΠ ο μεγαλύτερος.
\sol{7} α) Λάθος: το ΜΚΔ είναι 6. β) Λάθος: το ΕΚΠ είναι 12. γ) Σωστό. δ) Σωστό. ε) Σωστό. στ) Σωστό.
\sol{8} Το ΕΚΠ(15, 20): $15=3\cdot 5$, $20=2^2\cdot 5$, άρα $2^2\cdot 3\cdot 5=60$ δευτερόλεπτα. Στα 3 λεπτά (180 δευτ.) συμπίπτουν στα 60, 120 και 180 δευτερόλεπτα: 3 φορές.
\sol{9} $84=2^2\cdot 3\cdot 7$, $126=2\cdot 3^2\cdot 7$, άρα $\MKD=2\cdot 3\cdot 7=42$. Το πολύ 42 καλάθια, με $84:42=2$ μήλα και $126:42=3$ πορτοκάλια το καθένα.
\sol{10} $180=2^2\cdot 3^2\cdot 5$, $150=2\cdot 3\cdot 5^2$, άρα $\MKD=2\cdot 3\cdot 5=30$. Μέγιστη πλευρά 30 εκ. Πλακάκια: $180:30=6$ στο μήκος και $150:30=5$ στο πλάτος, δηλαδή $6\cdot 5=30$ πλακάκια.
\sol{11} $\EKP(15,\,20)=60$ λεπτά, άρα ξανά μαζί στις 8:00. Μαζί βρίσκονται στις 7:00, 8:00, 9:00 και 10:00: 4 φορές.
\sol{12} $\EKP(4,\,6,9)$: $4=2^2$, $6=2\cdot 3$, $9=3^2$, άρα $2^2\cdot 3^2=36$. Σε 36 μέρες.
\sol{13} $(12,\,18)$: ΜΚΔ 6, ΕΚΠ 36, $6\cdot 36=216=12\cdot 18$. $(8,\,20)$: 4 και 40, $4\cdot 40=160=8\cdot 20$. $(9,\,15)$: 3 και 45, $3\cdot 45=135=9\cdot 15$. $(14,\,21)$: 7 και 42, $7\cdot 42=294=14\cdot 21$. Το γινόμενο ΜΚΔ·ΕΚΠ είναι πάντα ίσο με το γινόμενο των δύο αριθμών.
\sol{14} $84=30\cdot 2+24$, και $\MKD(84,\,30)=6=\MKD(30,\,24)$. Συνεχίζω: $30=24\cdot 1+6$ και $24=6\cdot 4+0$. Ο τελευταίος διαιρέτης πριν το υπόλοιπο 0, το 6, είναι το ΜΚΔ. Για τα 252 και 105: $252=105\cdot 2+42$, $105=42\cdot 2+21$, $42=21\cdot 2+0$. Άρα $\MKD(252,\,105)=21$.
\sol{15} $8=2^3$ και $15=3\cdot 5$ δεν έχουν κοινό παράγοντα. Το ΕΚΠ πρέπει να περιέχει όλους τους παράγοντες και των δύο, με τον μεγαλύτερο εκθέτη, και επειδή δεν υπάρχει κανένας κοινός, τους παίρνει όλους: $2^3\cdot 3\cdot 5=8\cdot 15=120$.
\sol{16} Το γινόμενο των δύο αριθμών πρέπει να είναι $6\cdot 36=216$. Ένα ζευγάρι είναι το 6 και το 36 (ΜΚΔ 6, ΕΚΠ 36). Δεύτερο ζευγάρι: 12 και 18 (ΜΚΔ 6, ΕΚΠ 36, γινόμενο 216).
\sol{17} Όχι. Το ΜΚΔ είναι μικρότερο ή ίσο από τον μικρότερο αριθμό και το ΕΚΠ μεγαλύτερο ή ίσο από τον μεγαλύτερο. Αν οι αριθμοί είναι διαφορετικοί, ο μικρότερος είναι αυστηρά μικρότερος από τον μεγαλύτερο, άρα $\MKD<\EKP$. Ισότητα έχουμε μόνο όταν οι δύο αριθμοί είναι ίσοι.
